\section{Conclusion}\label{sec:conclusion}

The ProveIt iPRA development no longer primarily lacks semantic frameworks.  It has multiple mature
accounts of finite certificates, quantitative names, stable observers, demand closure, traces, and
residual producers.  The remaining general problem is compilation into the checked ordinary
calculus.

\TAPIR-iPRA addresses that problem with a deliberately small architecture:

\begin{enumerate}
\item specialize every fixed mathematical argument to concrete primitive-recursive syntax at host
      elaboration time;
\item represent the argument by an endpoint-indexed proof plan;
\item interpret the plan structurally through checked \(\Prf\) operations;
\item topologically order local proof nodes into a typed schedule;
\item compile earlier outputs away by simultaneous cut;
\item seal local and invariant certificates, and invoke bounded least search only for publicly named
      witnesses;
\item leave every genuinely unavailable producer as an explicit residual formula.
\end{enumerate}

The main mathematical theorem is small but decisive: a schedule whose local nodes are ordinary object
plans compiles to an ordinary object proof of every output from the residual context.  This closes the
last structural gap without internalizing a universal evaluator or inventing a second logic.

The current repository already checks the arithmetic mechanisms that make the architecture useful:
substitution, weakening, cut, quantifier-free induction, pairing, least search, bounded reflection,
bounded choice, and measure-decreasing termination.  The proposed new formal work is therefore
focused: elaborate the plan and schedule candidates, generalize bounded reflection to arbitrary
parameter arity, package reflected formulas compositionally, and migrate theorem packets one root at
a time.

For deep applications, the resulting theorem may remain conditional.  That is a strength.  A
statement such as
\[
  \mathsf{HeightBound}(q,H)\vdash_{\PRAAx}
  \mathsf{CompleteList}(q,H,L)
\]
identifies exactly what a future effective theorem must provide.  It separates the finite arithmetic
that iPRA can already verify from the domain theorem that has not yet been constructed.

The recommended next repository step is to integrate \code{Kernel.lean} and \code{Schedule.lean},
compile the three small examples listed in \cref{sec:lean}, and then use the schedule compiler to turn
one existing unverified packet into a literal residual sequent.  That experiment will measure the
true remaining cost of the last mile and provide a reusable pattern for the rest of the catalogue.
